% 本文件由 main.tex 输入，作为“实证结果分析”章节，不单独编译。

\section{实证结果分析}

\subsection{基准回归：政府产业引导基金与城市创新}

基准回归首先检验政府产业引导基金是否能够提升城市创新产出。表2报告了核心结果。可以看到，在控制城市固定效应、年份固定效应以及经济发展水平、财政科技支出、人口规模、产业结构和外资利用水平后，政府产业引导基金设立数量和累计设立规模的系数总体显著为正。其中，累计设立规模对发明申请量、实用新型申请量和专利申请总量均显著为正，说明基金规模扩张与城市专利申请产出提高存在稳定联系。

\begin{table}[htbp]
  \centering
  \caption{政府产业引导基金与城市创新：基准回归}
  {\fangsong\zihao{-5}
  \begin{tabularx}{0.96\linewidth}{l l c c c}
    \toprule
    因变量 & 基金口径 & 系数 & p值 & N \\
    \midrule
    发明申请量 & 累计设立规模 & 0.0300*** & $<0.001$ & 1083 \\
    实用新型申请量 & 累计设立规模 & 0.0405*** & 0.0031 & 1083 \\
    专利申请总量 & 累计设立规模 & 0.0800** & 0.0174 & 1083 \\
    发明申请量 & 当年设立数量 & 297.1917*** & $<0.001$ & 1149 \\
    专利申请总量 & 当年设立数量 & 849.5984*** & 0.0018 & 1149 \\
    发明申请量 & 累计设立数量 & 106.7968*** & $<0.001$ & 1149 \\
    \midrule
    控制变量 & \multicolumn{4}{c}{是} \\
    城市固定效应 & \multicolumn{4}{c}{是} \\
    年份固定效应 & \multicolumn{4}{c}{是} \\
    \bottomrule
  \end{tabularx}

  \par\smallskip
  {\songti\zihao{6}注：所有模型均按城市聚类稳健标准误。*、**、*** 分别表示在10\%、5\%、1\% 水平上显著。}
  }
\end{table}

% TODO：结果 CSV 解锁后，可按 README 要求为表2和表3补入各规格 $R^2$。

从经济含义看，基金累计设立规模对发明申请量和专利申请总量均显著为正，表明政府产业引导基金的作用并不限于增加一般性专利申请，也与高质量创新口径的改善相一致。当年设立数量和累计设立数量同样表现出显著正向结果，说明无论从基金数量还是基金规模观察，引导基金扩张都与城市创新产出提升相一致。由此，假说 H1 获得较强支持。考虑到累计设立规模在后续债务压力调节效应和机制检验中更为连贯，下文将其作为正文主解释变量。

\subsection{债务压力的调节效应}

在确认基金总体上具有创新促进效应后，本文进一步检验地方债务压力是否削弱这一作用。表3报告基金累计设立规模与滞后一期债务压力交互项的估计结果。交互项在专利申请总量、发明申请量和实用新型申请量三个口径下均显著为负，说明债务压力较高的城市中，同等规模政府产业引导基金所能带来的创新促进作用明显下降。

\begin{table}[htbp]
  \centering
  \caption{债务压力对基金创新效应的调节作用}
  {\fangsong\zihao{-5}
  \begin{tabularx}{0.92\linewidth}{l c c c}
    \toprule
    因变量 & 交互项系数 & p值 & N \\
    \midrule
    专利申请总量 & -0.0000609*** & $<0.001$ & 914 \\
    发明申请量 & -0.0000162*** & $<0.001$ & 914 \\
    实用新型申请量 & -0.0000292*** & $<0.001$ & 914 \\
    \midrule
    控制变量 & \multicolumn{3}{c}{是} \\
    城市固定效应 & \multicolumn{3}{c}{是} \\
    年份固定效应 & \multicolumn{3}{c}{是} \\
    \bottomrule
  \end{tabularx}

  \par\smallskip
  {\songti\zihao{6}注：交互项为“基金累计设立规模$\times$滞后一期债务压力”。所有模型均按城市聚类稳健标准误。同期债务压力口径下交互项同样显著为负。财政压力交互项在现有结果中方向与理论主线不一致，因此不纳入正文主结论。}
  }
\end{table}

该结果支持假说 H2。其含义是，政府产业引导基金并不是在任何财政环境下都能发挥同等作用。当地方债务压力上升时，地方政府的财政空间、风险承担能力和政策耐心可能下降，基金更容易受到短期绩效压力和偿债约束影响，从而削弱其对创新活动的边际促进作用。由于滞后一期债务压力更符合“债务约束先于基金绩效变化”的时间顺序，本文将表3作为调节效应的主规格。

\subsection{稳健性检验：替换变量口径}

为检验主结论是否依赖单一变量口径，本文从创新产出、基金规模和债务压力三个方面进行替换。首先，创新产出不仅使用专利申请总量，也使用发明申请量和实用新型申请量。基准回归中，基金累计设立规模对三类创新产出的系数均显著为正；调节效应中，基金累计设立规模与债务压力的交互项在三类创新产出下均显著为负。其次，基金变量除累计设立规模外，还使用当年设立数量和累计设立数量。结果显示，基金数量口径同样与城市创新产出显著正相关。最后，债务压力除滞后一期口径外，也使用当期债务压力和缩尾债务压力作为补充，核心交互项仍保持负向显著。

机制检验中的变量替换也提供了补充支持。耐心资本机制中，早期投资事件占比采用原始比例和1/99缩尾比例后均得到一致结果；金融发展机制中，存款/GDP 的反双曲正弦变换是正文主口径，对数变换口径方向相同但显著性略弱。综合来看，基准结论和债务压力调节结论并非由单一创新变量、单一基金变量或单一债务压力口径驱动。

\subsection{内生性处理：滞后变量检验}

本文进一步采用滞后变量缓解反向因果和同期冲击的担忧。债务压力调节效应的正文主规格使用滞后一期债务压力，而不是仅使用当期债务压力。表3显示，基金累计设立规模与滞后一期债务压力的交互项在三类创新产出下均显著为负，说明即便让债务压力先于创新结果进入模型，债务压力削弱基金创新效应的结论仍然成立。

机制检验也体现了创新产出的滞后响应。耐心资本机制中，早期投资事件占比主要影响未来两年专利申请总量；这与早期项目从投资到形成专利产出的时间滞后相一致。需要强调的是，滞后变量并不能完全消除所有内生性问题，尤其不能排除未观测的时变政策冲击。但在城市固定效应、年份固定效应和控制变量基础上，滞后债务压力与滞后创新响应结果增强了本文主结论的可信度。

\subsection{机制检验一：耐心资本属性}

第一条机制关注政府产业引导基金的耐心资本属性。理论上，债务压力越高，地方政府越可能要求基金更快见效、更低风险地配置资金，从而削弱其“投早、投小、投长期”的功能。本文使用早期投资事件占比代理耐心资本属性。为避免比例变量分母过小导致机械波动，机制检验将样本限定为基金投资事件数不少于2的城市—年份观测。该处理是识别比例机制的必要步骤，因为当分母为1时，单次投资会使早期投资占比机械地落在0或1，容易把偶然事件误判为结构性投资偏好。

\begin{table}[htbp]
  \centering
  \caption{耐心资本属性机制：早期投资事件占比}
  {\fangsong\zihao{-5}
  \resizebox{\linewidth}{!}{%
  \begin{tabular}{l l l c c c c c}
    \toprule
    创新结果 & 比例处理 & 债务处理 & 总效应 & 路径一 & 路径二 & $R^2$ & N \\
    \midrule
    两年后专利申请总量 & 原始比例 & 原始债务压力 & -0.085** & -0.143** & 0.092** & 0.680 & 266 \\
    两年后专利申请总量 & 比例1/99缩尾 & 原始债务压力 & -0.085** & -0.143** & 0.092** & 0.680 & 266 \\
    两年后专利申请总量 & 原始比例 & 债务压力1/99缩尾 & -0.075** & -0.128** & 0.091** & 0.680 & 266 \\
    两年后专利申请总量 & 比例1/99缩尾 & 债务压力1/99缩尾 & -0.075** & -0.128** & 0.091** & 0.680 & 266 \\
    \midrule
    控制变量 & \multicolumn{7}{c}{是} \\
    城市固定效应 & \multicolumn{7}{c}{是} \\
    年份固定效应 & \multicolumn{7}{c}{是} \\
    \bottomrule
  \end{tabular}%
  }

  \par\smallskip
  {\songti\zihao{6}注：总效应表示“基金规模$\times$债务压力”对未来创新的影响，路径一表示该交互项对早期投资事件占比的影响，路径二表示早期投资事件占比对未来创新的影响。表中结果来自基金投资事件数不少于2的分母稳定化样本。}
  }
\end{table}

表4显示，基金规模与债务压力的交互项显著降低早期投资事件占比，而早期投资事件占比对未来两年专利申请总量具有显著正向影响。与此同时，总效应显著为负。这说明债务压力可能通过削弱引导基金投向早期项目的耐心资本属性，进一步抑制其创新促进作用。需要谨慎的是，该组结果依赖分母稳定化样本，且早期投资金额占比口径并不稳定。因此，本文将其解释为具有透明样本边界的机制证据，而不将其写作无条件结论。

\subsection{机制检验二：社会资本撬动效率}

第二条机制关注社会资本撬动效率。政府产业引导基金的重要功能不仅在于直接提供财政资金，也在于通过政府背书和资本放大效应带动社会资本进入创新领域。债务压力上升时，地方政府可能降低基金运行稳定性、市场化程度和风险承担能力，从而削弱社会资本跟投和撬动效率。现有结果并不支持“社会资本撬动效率是普遍普通中介”的表述，而更支持一种调节式或承接机制：债务压力削弱撬动效率及其改善，而较高撬动效率会强化基金规模对创新的边际促进作用。

\begin{table}[htbp]
  \centering
  \caption{社会资本撬动效率机制：低金融发展地区}
  {\fangsong\zihao{-5}
  \resizebox{\linewidth}{!}{%
  \begin{tabular}{l l c c c c c}
    \toprule
    因变量 & 债务处理 & 基准交互项 & 机制方程 & 结果方程 & 加机制后交互项 & N \\
    \midrule
    实用新型申请量对数 & 滞后债务压力 & -4.03E-08*** & -8.95E-08** & 6.31E-07** & -3.68E-08*** & 705 \\
    实用新型申请量对数 & 滞后债务压力缩尾 & -4.03E-08*** & -8.95E-08** & 6.31E-07** & -3.68E-08*** & 705 \\
    下一期专利申请总量 & 滞后债务压力 & -8.20E-04** & -8.95E-08** & 0.02418** & -6.61E-04*** & 705 \\
    下一期专利申请总量 & 滞后债务压力缩尾 & -8.20E-04** & -8.95E-08** & 0.02418** & -6.61E-04*** & 705 \\
    \bottomrule
  \end{tabular}%
  }

  \par\smallskip
  {\songti\zihao{6}注：表中结果来自低金融发展、剔除2020年的样本，规格为未加控制变量的5\%全链条结果。机制变量为社会资本撬动效率年度变化，缺失补0后缩尾。加入机制项后，原债务调节项绝对值下降。}
  }
\end{table}

表5显示，在金融发展水平较低的城市中，债务压力会削弱政府引导基金撬动社会资本的效率或效率改善；而社会资本撬动效率越高，基金规模对创新产出的边际促进作用越强。加入社会资本撬动效率相关机制项后，原基金规模与债务压力交互项的绝对值下降，说明债务压力的抑制作用部分通过社会资本撬动效率这一承接渠道体现。由于最强证据集中在低金融发展地区，且5\%全链条结果来自未加控制规格，本文对该机制采用保守表述，将其限定为特定金融环境下更明显的调节式机制。

\subsection{机制检验三：融资约束缓解效应}

第三条机制检验政府产业引导基金能否通过改善地区金融发展环境来缓解融资约束。本文使用地区金融发展水平作为融资约束的反向代理，而不将其等同于企业层面的融资约束指数。在所有可用链条中，最适合正文报告的是存款/GDP 口径金融发展水平的反双曲正弦变换结果。

\begin{table}[htbp]
  \centering
  \caption{融资约束缓解机制：金融发展水平}
  {\fangsong\zihao{-5}
  \resizebox{\linewidth}{!}{%
  \begin{tabular}{l l c c c c c}
    \toprule
    机制变量口径 & 创新结果 & 交互项$\rightarrow$机制 & p值 & 机制$\rightarrow$创新 & p值 & N \\
    \midrule
    存款/GDP水平（asinh变换） & 专利申请总量对数 & -0.000000033 & 0.0147 & 0.4039 & 0.0497 & 944 \\
    贷款/GDP增速 & 实用新型申请量 & -0.001032 & 0.0228 & 10986.6163 & 0.0336 & 782 \\
    \bottomrule
  \end{tabular}%
  }

  \par\smallskip
  {\songti\zihao{6}注：第一行为正文主链条，表示基金累计设立规模与滞后债务压力交互项削弱金融发展水平，进而不利于创新产出。第二行为补充稳健性口径，其含义是金融发展改善速度。所有结果均纳入控制变量、城市固定效应和年份固定效应，并按城市聚类稳健标准误。}
  }
\end{table}

表6表明，当以存款/GDP 衡量金融发展水平并进行反双曲正弦变换时，基金累计设立规模与滞后债务压力的交互项对金融发展水平的影响显著为负，而金融发展水平对总专利申请对数显著为正。这说明债务压力会削弱政府产业引导基金改善地区金融环境、缓解融资约束的能力，并由此抑制城市创新产出。需要注意的是，主链条第二步的显著性接近5\%临界值，因此这一机制虽然与理论预期一致，但仍应以谨慎口径表述。

\subsection{本章小结}

综合来看，政府产业引导基金能够显著提升城市创新产出，但这一促进作用受到地方债务压力的制约。债务压力越高，基金累计设立规模对专利申请的边际促进效应越弱。进一步的机制检验表明，债务压力可能通过三条路径削弱引导基金的创新效应：一是降低基金投向早期项目的比例，削弱其耐心资本属性；二是在金融发展水平较低地区降低基金撬动社会资本的效率；三是削弱基金改善地区金融发展、缓解融资约束的能力。

上述结果也存在清晰边界。耐心资本机制依赖于基金投资事件数不少于2的分母稳定化样本；社会资本撬动效率机制主要出现在低金融发展地区，属于情境依赖较强的承接机制；金融发展机制虽然在理论链条上较为完整，但关键显著性处于5\%边界附近。因此，本文更适合将机制结果理解为与主结论方向一致、但具有样本条件和识别边界的补充证据。总体而言，政府产业引导基金的创新促进作用并非无条件成立，其政策绩效在很大程度上取决于地方财政债务约束与地区资本市场环境。
